% p006_a 的电脑重画（氢原子能级图：n=∞,5,4,3,2,1 及能量；三条跃迁箭头 hν_1: n=1→4，hν_2: n=2→4，hν_3: n=1→2）
\begin{tikzpicture}[line width=0.7pt, >={Stealth[length=2.2mm]}]
  % 能级线
  \draw (0.0,3.05) -- (2.03,3.05);
  \draw (-0.07,2.77) -- (2.0,2.77);
  \draw (-0.05,2.21) -- (1.98,2.21);
  \draw (-0.07,1.71) -- (2.0,1.71);
  \draw (-0.13,0.99) -- (2.07,0.99);
  \draw (-0.23,0.0) -- (2.25,0.0);
  % 左侧 n 标注
  \node at (-1.3,3.15) {$n=\infty$};
  \node at (-1.53,2.63) {$n=5$};
  \node at (-1.38,2.15) {$n=4$};
  \node at (-1.48,1.7) {$n=3$};
  \node at (-1.4,0.93) {$n=2$};
  \node at (-1.08,-0.07) {$n=1$};
  % 右侧能量标注
  \node[right] at (2.03,3.05) {$0$};
  \node[right] at (2.3,2.8) {$-0.54$};
  \node[right] at (2.4,2.25) {$-0.85$};
  \node[right] at (2.45,1.75) {$-1.5\unit{eV}$};
  \node[right] at (2.65,0.95) {$-3.4\unit{eV}$};
  \node[right] at (2.55,0.02) {$-13.6\unit{eV}$};
  % 跃迁箭头
  \draw[->] (0.21,-0.22) -- (0.21,2.17);
  \draw[->] (1.1,0.99) -- (1.1,2.21);
  \draw[->] (1.0,0.0) -- (1.0,0.99);
  \node at (-0.2,0.55) {$h\nu_1$};
  \node at (1.52,1.35) {$h\nu_2$};
  \node at (1.5,0.33) {$h\nu_3$};
\end{tikzpicture}
